\documentclass{beamer}

\usetheme{metropolis} 
\usepackage[utf8]{inputenc}
\usepackage{graphicx}
\usepackage{tabularx}
\usepackage{booktabs}
\usepackage{array} 
\newcolumntype{Y}{>{\raggedright\arraybackslash}X} 
\usepackage{caption} 
% font=scriptsize makes the text smaller (you can also use 'tiny') 
% skip=2pt reduces the vertical space between the image and caption 
\captionsetup{font=scriptsize, labelfont=scriptsize, skip=2pt}

\title{Classifying Parliamentary Text into Ministries}
\subtitle{A Study of Representation Methods \\ [1em] DS5601: Natural Language Processing}
\author{142502012 --- Devadatta Mahesh Pokharanakar \\ 142502010 --- Chilaka Sri Krishna Sai \\ 142502032 --- Vaddi Govardhan}
\institute{Department of Data Science \\ IIT Palakkad}
\date{\today}

\begin{document}

% Title Slide
\maketitle

% Slide 1
\begin{frame}{Problem Motivation}
\begin{itemize}
    \item Parliamentary proceedings are released as long, unstructured documents.
    \item Important information (topics, ministries) is not explicitly labeled.
    
    \item \textbf{Goal:} Identify which ministry a paragraph belongs to.
    
    \item \textbf{Key Challenge:}
    \begin{itemize}
        \item No annotated dataset
        \item High domain complexity
    \end{itemize}
    
    \item \textbf{Core Idea:} 
    \begin{itemize}
        \item Explore different paradigms of text representation
        \item Evaluate methods without relying on gold labels
    \end{itemize}
\end{itemize}
\end{frame}

% Slide 2
\begin{frame}{Data and Setup}
\begin{itemize}
    \item \textbf{Data:}
    \begin{itemize}
        \item Parliamentary transcripts (debates, budget speeches, Q\&A)
        \item 53 ministry knowledge bases (Wikipedia + Govt + LLM summaries)
    \end{itemize}
    
    \item \textbf{Processing:}
    \begin{itemize}
        \item Speaker-wise and paragraph-level segmentation
        \item Cleaning + domain-specific stopwords
    \end{itemize}
    
    \item \textbf{Evaluation:}
    \begin{itemize}
        \item No ground truth labels
        \item Weak supervision using LLM + cross-method comparison
    \end{itemize}
\end{itemize}
\end{frame}

% Slide 3
\begin{frame}{Topic Modeling (Supporting Structure)}
\begin{itemize}
    \item Topic modeling performed using \textbf{David Mimno’s online tool} \small \url{https://mimno.infosci.cornell.edu/jsLDA/}
    
    \item \textbf{Outputs obtained:}
    \begin{itemize}
        \item 25 topics over the full corpus
        \item Top 10 representative words per topic
        \item Document--Word--Topic state distributions
        \item Automatically generated domain-specific stopwords
    \end{itemize}
\end{itemize}

\begin{figure}
    \centering
    \includegraphics[width=0.9\textwidth]{images/word-topic-doc.png}
    \vspace{-0.5em}
    \caption{Example of word--topic and word--document relationships}
\end{figure}
\end{frame}

% Slide 4
\begin{frame}{Three Ways to Represent Text}
\centering
\begin{tabular}{@{} l l @{}}
\toprule
\textbf{Method} & \textbf{Core Idea} \\
\midrule
Lexical (TF-IDF) & Meaning = Words \\
Probabilistic & Words $\rightarrow$ Topics $\rightarrow$ Ministries \\
Embedding-based & Meaning = Vector similarity \\
\bottomrule
\end{tabular}

\vspace{1em}

\textbf{Research Question:}  
Which representation works best when no labels are available?
\end{frame}

% Slide 5
\begin{frame}{Method 1: Lexical (TF-IDF)}
\textbf{Idea:} Match paragraph words with ministry-specific vocabulary

\vspace{0.5em}

\textbf{Assumption:}
\begin{itemize}
    \item Words directly indicate meaning
\end{itemize}

\textbf{Works well when:}
\begin{itemize}
    \item Distinct domain keywords exist (e.g., ``fertilizer'', ``crop'')
\end{itemize}

\textbf{Fails when:}
\begin{itemize}
    \item Synonyms are used (``budget'' vs ``fiscal policy'')
    \item Vocabulary overlaps across ministries
\end{itemize}

\vspace{0.5em}
\textbf{Insight:} Parliamentary language is not discriminative enough for pure keyword methods.
\end{frame}

% Slide 6
\begin{frame}{Method 2: Probabilistic (Doc-Word-Topic)}
\textbf{Idea:} Words $\rightarrow$ Topics $\rightarrow$ Ministries

$Score(d) = \sum_{w} \sum_{t} P(d|w)\cdot P(t|w)\cdot P(d|t)$

\textbf{Assumptions:}
\begin{itemize}
    \item Words are independent
    \item Topics explain word usage
    \item Frequencies approximate probabilities
\end{itemize}

\textbf{Strength:} Combines corpus statistics with structure

\textbf{Fails because:}
\begin{itemize}
    \item Independence assumption is violated
    \item Topics overlap across ministries
\end{itemize}

\textbf{Insight:} Real language does not follow clean probabilistic assumptions.
\end{frame}

% Slide 7
\begin{frame}{Method 3: Embedding-Based}
\textbf{Idea:} Represent text as vectors and compare using similarity

\textbf{Variants:}
\begin{itemize}
    \item Whole-document embedding (loses detail)
    \item Mean of paragraph embeddings (\textbf{best})
    \item Top-K similar paragraphs (captures local similarity)
\end{itemize}

\textbf{Why it works:}
\begin{itemize}
    \item Captures semantic meaning beyond keywords
    \item Handles synonyms and paraphrases
\end{itemize}

\textbf{Limitations:}
\begin{itemize}
    \item Sensitive to representation design
    \item Can blur multiple topics
\end{itemize}
\end{frame}

% Slide 8
\begin{frame}{Embedding Space Visualization}
\begin{columns}
    \begin{column}{0.5\textwidth}
        \begin{figure}
            \centering
            \includegraphics[width=\textwidth]{images/tsne-embed.png}
            \caption{t-SNE projection}
        \end{figure}
    \end{column}
    
    \begin{column}{0.5\textwidth}
        \begin{figure}
            \centering
            \includegraphics[width=\textwidth]{images/pca-embed.png}
            \caption{PCA projection}
        \end{figure}
    \end{column}
\end{columns}

\textbf{Observations:}
\begin{itemize}
    \item Some ministries form clear clusters
    \item Significant overlap exists between related ministries
\end{itemize}
\end{frame}

% Slide 9
\begin{frame}{Clustering Ministries}
\begin{itemize}
    \item Observed weak separation (low silhouette scores) for hard clustering
    \item Adopted \textbf{consensus-based clustering}: Group ministries based on stable co-occurrence patterns
\end{itemize}

\begin{figure}
    \centering
    \includegraphics[width=0.85\textwidth]{images/cluster-formation.png}
    \caption{\scriptsize Consensus-based grouping of ministries}
\end{figure}

\vspace{0.3em}

\textbf{Insight:} Ministries form soft clusters rather than clearly separable classes.
\end{frame}

% Slide 10
\begin{frame}{Comparison with Weak Supervision (ChatGPT Labels)}
\vspace{0.5em}
\begin{columns}
    \begin{column}{0.33\textwidth}
        \begin{figure}
            \centering
            \includegraphics[width=\textwidth]{images/exact-match.png}
            \caption{Exact Match}
        \end{figure}
    \end{column}
    
    \begin{column}{0.33\textwidth}
        \begin{figure}
            \centering
            \includegraphics[width=\textwidth]{images/topk-errors.png}
            \caption{Top-K Errors}
        \end{figure}
    \end{column}
    
    \begin{column}{0.33\textwidth}
        \begin{figure}
            \centering
            \includegraphics[width=\textwidth]{images/cluster-match.png}
            \caption{Cluster Match}
        \end{figure}
    \end{column}
\end{columns}

\textbf{Observations:}
\begin{itemize}
    \item Embedding methods show higher consistency across metrics
    \item Probabilistic method improves over lexical baseline
    \item Exact match is low, but cluster-level agreement is higher
\end{itemize}

\textbf{Insight:} In absence of ground truth, consistency across views is more meaningful than raw accuracy.
\end{frame}

% Slide 11
\begin{frame}{Results and Key Observations}
\begin{itemize}
    \item Embedding-based methods outperform lexical approaches
    \item Probabilistic method improves over simple keyword matching
    \item Best performance: \textbf{paragraph-level embedding averaging}
\end{itemize}

\vspace{0.5em}

\textbf{Important Observation:}
\begin{itemize}
    \item No method is perfectly correct
    \item Some methods are more \textbf{consistent} than others
\end{itemize}

\vspace{0.5em}

\textbf{Reason:}
\begin{itemize}
    \item High semantic overlap between ministries
\end{itemize}
\end{frame}

% Slide 12
\begin{frame}{Conclusion and Takeaways}
\begin{itemize}
    \item Ministry classification is fundamentally a \textbf{representation problem}
    
    \item \textbf{Lexical methods:}
    \begin{itemize}
        \item Fail due to vocabulary overlap and lack of context
    \end{itemize}
    
    \item \textbf{Probabilistic methods:}
    \begin{itemize}
        \item Limited by unrealistic independence and topic assumptions
    \end{itemize}
    
    \item \textbf{Embedding methods:}
    \begin{itemize}
        \item Capture semantics effectively
        \item Performance depends on representation design (granularity)
    \end{itemize}
\end{itemize}

\vspace{1em}

\textbf{Key Insight:} \centering Better models did not solve the problem \\ --- better representations did.
\end{frame}

% Slide 13
\begin{frame}
\centering
\Huge Thank You

\vspace{1em}

\Large Questions?

\vfill

\scriptsize Project Repository: \\ \url{https://github.com/DevadattaP/analyse_parliament_transcripts_NLP}
\end{frame}

\end{document}
